\documentclass[krantz1]{krantz}
\usepackage{fixltx2e,fix-cm}
\usepackage{amssymb, bbm, bm, mathtools}
\usepackage{amsmath}
\usepackage[pdftex]{graphicx}
\usepackage{color, rotating}
\usepackage{subcaption}
\usepackage{hyperref}
\usepackage{cleveref}
\usepackage{xeCJK}
\usepackage{natbib}

% Theorem-like environments
\newtheorem{exercise}{Exercise}[chapter]
\newtheorem{example}{Example}
\newtheorem{definition}{Definition}
\newtheorem{assumption}{Assumption}[chapter]
\newtheorem{thesis}{Theorem}[chapter]
\newtheorem{corollary}{Corollary}[chapter]

\title{Problem Set 5: Comprehensive Solutions}
\author{乐绎华 学号 23363017}

% Custom environments from the book template
\newenvironment{myproof}[2] {\noindent {\bfseries Proof of {#1} {#2}:}}{\hfill$\square$}
\newenvironment{mysolution}[1]{\vspace{0.1in} \noindent {\bfseries Solution to Problem {#1}:}}{\hfill$\square$}

% Custom macros from the book
\def\iidsim{\stackrel{\textup{IID}}{\sim}}
\def\indsim{\stackrel{ind}{\sim}}
\def\convergeas{\stackrel{a.s.}{\longrightarrow}}
\def\converged{\stackrel{\textup{d}}{\longrightarrow}}
\def\pr{\textup{pr}}
\def\cov{\textup{cov}}
\def\ind{\perp}
\def\RD{\textsf{rd}}
\def\RR{\textsf{rr}}
\def\OR{\textsf{or}}
\def\true{\textup{true}}
\def\obs{\textup{obs}}
\def\se{\textup{se}}
\def\logit{\textup{logit}}
\def\expit{\textup{expit}}
\def\sumn{\sum_{i=1}^n}
\def\diff{\textsf{d}}
\def\Paragraph{\paragraph}
\def\at{\textup{a}}
\def\cp{\textup{c}}
\def\df{\textup{d}}
\def\nt{\textup{n}}
\def\CACE{\text{CACE}}
\def\TSLS{\text{TSLS}}
\def\ILS{\text{ILS}}
\def\var{\textup{var}}
\def\sumn{\sum_{i=1}^n}
\def\tran{\textup{T}}

\begin{document}
\maketitle

\chapter{Problem Set 5}

This problem set covers Chapters 21, 22, and 23 of Peng Ding's \textit{A First Course in Causal Inference}, focusing on the Instrumental Variable (IV) method from both experimental and econometric perspectives.

\section{Background: The Instrumental Variable Framework}

Consider units $i=1, \dots, n$. Let $Z_i$ be the treatment assignment (instrument), $D_i$ be the treatment received, and $Y_i$ be the outcome. 

\begin{definition}[Latent Compliance Status]\label{def::compliance}
Stratify the population into four latent groups based on $U_i = \{D_i(1), D_i(0)\}$:
\begin{itemize}
    \item Always-taker ($\at$): $D_i(1)=1, D_i(0)=1$.
    \item Complier ($\cp$): $D_i(1)=1, D_i(0)=0$.
    \item Defier ($\df$): $D_i(1)=0, D_i(0)=1$.
    \item Never-taker ($\nt$): $D_i(1)=0, D_i(0)=0$.
\end{itemize}
\end{definition}

\begin{assumption}[randomization]\label{21.1}
$Z \ind \{D(1), D(0), Y(1), Y(0)\}$.
\end{assumption}

\begin{assumption}[monotonicity]\label{21.2}
$\pr(U = \df) = 0$, i.e., $D_i(1) \ge D_i(0)$ for all $i$.
\end{assumption}

\begin{assumption}[exclusion restriction]\label{21.3}
$Y_i(1) = Y_i(0)$ for always-takers ($U=\at$) and never-takers ($U=\nt$).
\end{assumption}

\begin{thesis}\label{thm::21.1}\emph{(Theorem 21.1, Ding, 2025)}
Under Assumptions \ref{21.1}--\ref{21.3}, the Complier Average Causal Effect (CACE) is identified as:
$$ \tau_\cp = E\{Y(1) - Y(0) \mid U=\cp\} = \frac{E(Y\mid Z=1) - E(Y\mid Z=0)}{E(D\mid Z=1) - E(D\mid Z=0)}. $$
\end{thesis}

\section{Exercises from Chapter 21: Compliance, IV, and LATE}

\Paragraph{21.1~Variance of the Wald estimator}\label{hw::infinite-var-iv}

Show that $\var(\hat{\tau}_\cp)=\infty$.

\begin{mysolution}{21.1}
The Wald estimator is $\hat{\tau}_\cp = \hat{\tau}_Y / \hat{\tau}_D$. In finite samples, $\hat{\tau}_D$ is a discrete random variable with a strictly positive probability that $\hat{\tau}_D = 0$ (e.g., if all units in the sample are never-takers). Consequently, the expectation of $\hat{\tau}_\cp^2$ diverges due to the non-zero probability of a zero denominator, leading to $\var(\hat{\tau}_\cp) = \infty$.
\end{mysolution}

\Paragraph{21.6~Binary IV and ordinal treatment received}\label{para::iv-ordinal-treatment}

\citet{angrist1995two} discussed a binary IV $Z$, an ordinal treatment $D \in \{0, 1, \dots, J\}$, and outcome $Y$. Under randomization, monotonicity ($D(1) \ge D(0)$), and exclusion restriction ($Y(z,d) = Y(d)$), prove:
$$ \frac{E(Y\mid Z=1) - E(Y\mid Z=0)}{E(D\mid Z=1) - E(D\mid Z=0)} = \sum_{j=1}^J w_j E\{Y(j) - Y(j-1) \mid D(1) \ge j > D(0)\} $$
where $w_j = \frac{\pr\{D(1) \ge j > D(0)\}}{\sum_{j'=1}^J \pr\{D(1) \ge j' > D(0)\}}$.

\begin{mysolution}{21.6}
Write $D(z) = \sum_{j=1}^J \mathbb{I}(D(z) \ge j)$ and $Y(D(z)) = Y(0) + \sum_{j=1}^J [Y(j) - Y(j-1)] \mathbb{I}(D(z) \ge j)$.
By monotonicity, $\mathbb{I}(D(1) \ge j) - \mathbb{I}(D(0) \ge j) = \mathbb{I}(D(1) \ge j > D(0))$.
The denominator of the Wald ratio is:
$$ E[D(1) - D(0)] = \sum_{j=1}^J \pr(D(1) \ge j > D(0)). $$
The numerator (ITT effect on $Y$) is:
\begin{align*}
E[Y(D(1)) - Y(D(0))] &= \sum_{j=1}^J E\{ [Y(j) - Y(j-1)] \mathbb{I}(D(1) \ge j > D(0)) \} \\
&= \sum_{j=1}^J E\{ Y(j) - Y(j-1) \mid D(1) \ge j > D(0) \} \pr(D(1) \ge j > D(0)).
\end{align*}
Dividing the numerator by the denominator yields the weighted average with weights $w_j$.
\end{mysolution}

\Paragraph{21.7~Data analysis: a flu shot encouragement design}\label{hw::flushot}

Analyze \texttt{fludata.txt} with and without covariates.

\begin{mysolution}{21.7}
The analysis yields the following estimates:
\begin{itemize}
    \item \textbf{Unadjusted}: $\hat{\tau}_Y = -0.0147$, $\hat{\tau}_D = 0.1184 \implies \hat{\tau}_\cp = \mathbf{-0.1246}$.
    \item \textbf{Lin-adjusted}: $\hat{\tau}_{Y,L} = -0.0149$, $\hat{\tau}_{D,L} = 0.1187 \implies \hat{\tau}_{\cp,L} = \mathbf{-0.1252}$.
\end{itemize}
Encouraging flu shots reduces hospitalization by $\sim$12.5\% among compliers.
\end{mysolution}

\section{Exercises from Chapter 22: IV Inequalities}

\begin{thesis}\label{thm::22.1}\emph{(Theorem 22.1, Ding, 2025)}
Under Assumptions \ref{21.1}--\ref{21.3}:
$$ \mu_{1\cp} = \frac{E(DY\mid Z=1) - E(DY\mid Z=0)}{\pi_\cp}, \quad \mu_{0\cp} = \frac{E\{(1-D)Y\mid Z=0\} - E\{(1-D)Y\mid Z=1\}}{\pi_\cp} $$
where $\pi_\cp = E(D\mid Z=1) - E(D\mid Z=0)$.
\end{thesis}

\Paragraph{22.2~Risk ratio for compliers}\label{hw::rr-compliers}

Define $\RR_\cp = \frac{\pr\{Y(1)=1\mid U=\cp\}}{\pr\{Y(0)=1\mid U=\cp\}}$. Show:
$$ \RR_\cp = \frac{E(DY\mid Z=1) - E(DY\mid Z=0)}{E\{(D-1)Y\mid Z=1\} - E\{(D-1)Y\mid Z=0\}}. $$

\begin{mysolution}{22.2}
By Theorem \ref{thm::22.1}, $\mu_{1\cp} = \pi_\cp^{-1}[E(DY\mid Z=1) - E(DY\mid Z=0)]$ and $\mu_{0\cp} = \pi_\cp^{-1}[E((1-D)Y\mid Z=0) - E((1-D)Y\mid Z=1)]$.
Observe $E\{(D-1)Y \mid Z\} = E\{DY - Y \mid Z\} = -E\{(1-D)Y \mid Z\}$.
Thus, $E\{(D-1)Y \mid Z=1\} - E\{(D-1)Y \mid Z=0\} = E\{(1-D)Y \mid Z=0\} - E\{(1-D)Y \mid Z=1\}$.
The ratio $\mu_{1\cp}/\mu_{0\cp}$ then simplifies to the target identity.
\end{mysolution}

\Paragraph{22.7~Bounds on the average causal effect on the whole population}\label{para::bounds-ACE}

For $Y \in [\underline{y}, \overline{y}]$, let $\delta = E\{Y(d=1) - Y(d=0)\}$. Prove:
$$ \underline{\delta} = \delta' - \overline{y} \pr(D=1\mid Z=0) + \underline{y} \pr(D=0\mid Z=1) $$
$$ \overline{\delta} = \delta' - \underline{y} \pr(D=1\mid Z=0) + \overline{y} \pr(D=0\mid Z=1) $$
where $\delta' = E(DY\mid Z=1) - E((1-D)Y\mid Z=0)$.

\begin{mysolution}{22.7}
Decompose $\delta = \pi_\at(m_{1\at} - m_{0\at}) + \pi_\nt(m_{1\nt} - m_{0\nt}) + \pi_\cp(m_{1\cp} - m_{0\cp})$.
Define $\delta' = \pi_\at m_{1\at} + \pi_\cp (m_{1\cp} - m_{0\cp}) - \pi_\nt m_{0\nt}$.
Then $\delta = \delta' + \pi_\nt m_{1\nt} - \pi_\at m_{0\at}$.
Since $m_{1\nt}, m_{0\at} \in [\underline{y}, \overline{y}]$, the lower bound occurs at $(m_{1\nt}, m_{0\at}) = (\underline{y}, \overline{y})$ and the upper bound at $(\overline{y}, \underline{y})$. Substitution yields the sharp bounds.
\end{mysolution}

\section{Exercises from Chapter 23: Econometric Perspective and TSLS}

\Paragraph{23.1~More algebra for TSLS}\label{hw::algebra-tsls}
\begin{enumerate}
    \item Show $\hat{\Gamma} = (\sumn Z_i Z_i\tran)^{-1} \sumn Z_i D_i\tran$.
    \item Show $\hat{\beta}_\TSLS = \hat{\beta}_{\textsc{iv}}$ in the just-identified case.
\end{enumerate}

\begin{mysolution}{23.1}
1. The OLS estimator $\hat{\Gamma}$ minimizes $\sumn \|D_i - \Gamma^\tran Z_i\|^2$, satisfying the normal equation $\sumn Z_i(D_i - \hat{\Gamma}^\tran Z_i)^\tran = 0$. Solving for $\hat{\Gamma}$ yields $\hat{\Gamma} = (\sumn Z_i Z_i\tran)^{-1} \sumn Z_i D_i\tran$.

2. If $Z, D$ have the same dimension, $\hat{\Gamma}$ is invertible.
$\sum \hat{D}_i \hat{D}_i\tran = \hat{\Gamma}^\tran (\sum Z_i Z_i\tran) \hat{\Gamma} = (\sum D_i Z_i\tran) (\sum Z_i Z_i\tran)^{-1} (\sum Z_i D_i\tran)$.
Similarly, $\sum \hat{D}_i Y_i = \hat{\Gamma}^\tran \sum Z_i Y_i = (\sum D_i Z_i\tran) (\sum Z_i Z_i\tran)^{-1} \sum Z_i Y_i$.
Taking the ratio cancels $(\sum D_i Z_i\tran)$ and $(\sum Z_i Z_i\tran)^{-1}$, leaving $\hat{\beta}_\TSLS = (\sum Z_i D_i\tran)^{-1} \sum Z_i Y_i = \hat{\beta}_{\textsc{iv}}$.
\end{mysolution}

\Paragraph{23.2~Equivalence between TSLS and ILS}\label{para::tsls-ils}
Prove $\hat{\beta}_{1,\ILS} = \hat{\beta}_{1,\TSLS}$ for a single treatment and IV.

\begin{mysolution}{23.2}
The ILS estimator is $\hat{\beta}_{1,\ILS} = \hat{\Gamma}_1 / \hat{\gamma}_1 = \frac{\cov(Z, Y)/\var(Z)}{\cov(Z, D)/\var(Z)} = \frac{\cov(Z, Y)}{\cov(Z, D)}$.
TSLS regresses $Y$ on $\hat{D} = \hat{\gamma}_0 + \hat{\gamma}_1 Z$:
$\hat{\beta}_{1,\TSLS} = \frac{\cov(\hat{D}, Y)}{\var(\hat{D})} = \frac{\cov(\hat{\gamma}_1 Z, Y)}{\var(\hat{\gamma}_1 Z)} = \frac{\hat{\gamma}_1 \cov(Z, Y)}{\hat{\gamma}_1^2 \var(Z)} = \frac{\cov(Z, Y)}{\hat{\gamma}_1 \var(Z)}$.
Substituting $\hat{\gamma}_1 = \cov(Z, D)/\var(Z)$ yields $\frac{\cov(Z, Y)}{\cov(Z, D)}$.
\end{mysolution}

\Paragraph{23.3~Control function in the linear IV model}\label{para::control-function-iv}
Show $\hat{\beta}_{\textsc{cf}} = \hat{\beta}_\TSLS$.

\begin{mysolution}{23.3}
By the FWL theorem, the coefficient of $D$ in a regression of $Y$ on $(D, \check{D})$ is the same as regressing $Y$ on the residual of $D$ on $\check{D}$.
Since $D = \hat{D} + \check{D}$ and $\hat{D} \perp \check{D}$, this residual is exactly $\hat{D}$. Regressing $Y$ on $\hat{D}$ is the TSLS procedure, hence the equivalence.
\end{mysolution}

\section{Additional Problems}

\Paragraph{2.~CACE as Ratio of Covariances}\label{hw::wald-cov-form}
Show $\CACE = \frac{\cov(Z, Y)}{\cov(Z, D)}$.

\begin{mysolution}{2}
For binary $Z$, $\cov(Z, Y) = \pr(Z=1)\pr(Z=0) [E(Y\mid Z=1) - E(Y\mid Z=0)]$ and $\cov(Z, D) = \pr(Z=1)\pr(Z=0) [E(D\mid Z=1) - E(D\mid Z=0)]$. Their ratio is the Wald identity, which identifies CACE.
\end{mysolution}

\Paragraph{3.~Weighting method for CACE in observational studies}\label{hw::weighting-cace-x}
Show $E[\kappa_{1i} g(Y, X)] = \pr(U=\cp) E[g(Y(1), X) \mid U=\cp]$.

\begin{mysolution}{3}
By iterated expectations, $E[\kappa_{1i} g(Y, X)] = E_X [ E[D \frac{Z - e(X)}{e(X)(1-e(X))} g(Y, X) \mid X ] ]$.
The inner term is $e(X)(1-e(X))^{-1} [E(DZ g \mid X) - e(X) E(D g \mid X)]$.
Simplifying, $E(DZ g \mid X) = e(X) E[D(1) g(1) \mid X]$ and $E(D g \mid X) = e(X) E[D(1) g(1) \mid X] + (1-e(X)) E[D(0) g(0) \mid X]$.
Substitution yields $E_X[E((D(1)-D(0))g(1) \mid X)] = E[\mathbb{I}(U=\cp)g(1)] = \pr(U=\cp) E[g(1) \mid U=\cp]$.
\end{mysolution}

\Paragraph{4.~Analysis of Card (1993)}\label{hw::card-schooling}
Discuss dichotomization vs. TSLS.

\begin{mysolution}{4}
Dichotomization identifies LATE for a specific threshold, providing clear interpretability for a sub-population. TSLS utilizes the full education distribution and, under linearity, identifies the average derivative effect.
\end{mysolution}

\bibliographystyle{plainnat}
\bibliography{causal/causal}

\end{document}
